\documentclass[10pt,a4paper]{amsart}
\usepackage[margin=19mm]{geometry}
\usepackage{amsmath,amssymb,mathtools}
\usepackage[scr=rsfs,cal=euler]{mathalfa}
\usepackage{xfrac}
\usepackage[unicode,hidelinks,bookmarks=false]{hyperref}
\newcommand{\ie}{i.e.\ }
\newcommand{\cf}{cf.\ }
\newcommand{\resp}{respectively\ }
\newcommand{\Subsection}{\subsection}
\providecommand{\nobreakdash}{\nobreak\hbox{-}}
\setlength{\emergencystretch}{2em}
\multlinegap0pt
\usepackage{amssymb}
\usepackage{tikz}
\usepackage{float}
\newtheorem{proposition}{Proposition}
\def\FPP {{\P^2_{fake}}}
\def\P {{\mathbb P}}
\def\Z {{\mathbb Z}}
\title{Finding equations of the fake projective plane $(C18,p=3,\{2I\})$}
\author{Lev Borisov, Bojue Wang}
\date{Source: arXiv:2512.03213v2 (CC BY 4.0)}
\begin{document}
\maketitle

\noindent\emph{Source and licence.} Finding equations of the fake projective plane $(C18,p=3,\{2I\})$. Version arXiv:2512.03213v2 (CC BY 4.0), distributed by arXiv under the Creative Commons Attribution 4.0 International licence, http://creativecommons.org/licenses/by/4.0/, as recorded in the arXiv licence metadata for this exact version and shown on the abstract page https://arxiv.org/abs/2512.03213v2. Full licence notice: \texttt{licenses/arxiv-2512.03213-CC-BY-4.0.txt}.\par\smallskip

\noindent\textit{Editorial scope.} Counted unit: the article's proposition key (source0.tex lines 472--490). The article's complete original proof of it is delivered with it (source0.tex lines 492--498, one \texttt{proof} environment of 7 lines). No mathematics is written, repaired or re-derived: the statement and the proof are the article's own lines, in the article's own notation, and no retained line is substituted or re-flowed.  No further source-local context is needed: every cross-reference of every retained line resolves inside the two delivered spans.  Nothing was omitted from any retained span, so the package carries no omission record.  No retained line contains a citation, so the wrapper carries no bibliography and no bibliography entry.  No retained line contains a figure environment, an image-inclusion command or drawing code, so no image is delivered and the package declares no asset dependency.  Editorial formatting only, in addition: an AMS \texttt{amsart} preamble that restates the article's own declarations and macros used by the retained lines, namely the environments \texttt{proposition} on separate counters, together with the standard packages those lines need.  The retained environments print in this document as Proposition 1 in order of appearance, whereas the article prints the same items with the numbering of its own full document.  The complete native source of the article as distributed in the arXiv e-print for this exact version is bundled unchanged as \texttt{sources/190244/source0.tex}.
\begin{proposition}\label{key}
Consider an order $3$ element $h$ of $SL(2,\Z/3\Z)$ and its action on the $3$-torsion  subgroup $C_3\times C_3$ of ${\rm Pic}(8.\FPP)$. Suppose that the 
character $(w,1)$ of the covering $C_3\times C_3$ corresponds to the eigenvector of $h$ in $C_3\times C_3$.
Let $f_1$ and $f_2 = \sigma(f_1)$ be a $(w,1)$-eigenvector and a  $(w^2,1)$-eigenvector for the covering $C_3\times C_3$ in the space $H^0(72.\FPP,3H)_{29}$, respectively.
Likewise, let $g_1$ and $g_2=\sigma(g_1)$ be an $(w,1)$- and $(w^2,1)$-eigenvectors  in the space $H^0(72.\FPP,3H)_{30}$. Then  $s_1 = f_1 f_2$, $s_2 = f_1 g_2$, $s_3 = g_1 f_2$, $s_4 = g_1 g_2 $
are invariant under the covering group and can be thought of as elements of $H^0(8.\FPP,6H)$. These sections $s_i$ have the following properties.
\begin{itemize}
\item
$
s_1s_4 = s_2 s_3
$ 
\item 
$\sigma(s_2) = s_3$
\item
Sections $s_1,s_2, s_3,s_4$ have weights $(w,1,1,w^2)$ respectively for the central $C_3$ action on $H^0(8.\FPP,6H)$. 
\item
The weights of $s_1,s_2,s_3,s_4$ for the action of $h\in SL(2,\Z/3\Z)$ are $(w^{2a},w^{a+b},w^{a+b},w^{2b})$ for some $a$ and $b$ in $\Z/3\Z$.
\end{itemize}
\end{proposition}

\begin{proof}
The first two statements are immediate from the construction. To prove the third statement, observe that the generator of the central $C_3$  has trace $8w^2$
in $\chi_{29}$ and $8w$ in $\chi_{30}$ (after an appropriate choice of generator or a switch of $\chi_{29}$ and $\chi_{30}$). Thus $f_i$ have eigenvalues $w^2$ and $g_i$ have eigenvalues $w$. 

\medskip
The last statement is the most delicate. Since $h$ preserves the corresponding element of the Picard group, its action preserves the corresponding eigenspaces of $H^0(72.\FPP,3H)_{29}$ and $H^0(72.\FPP,3H)_{30}$. Thus $f_i$ and $g_i$ are eigenvectors for its action, with eigenvalues $w^a$ and $w^b$ for some $a$ and $b$.
\end{proof}

\end{document}
